\subsection{Discussion}

The advantage of the Jev-like organization grows with the number of pockets. For 57 pockets the featurization of the ligands, which every network needs, dominates the time of the dual encoder, and the advantage over the cross-attention network is small. The organization therefore pays in the screen of a library against many pockets, where all pockets share the ligand embedding.

The price of the organization is a few points of recall at a budget of the best 10\% of the ranking. The dual encoder keeps \kfRatioDual{}\% of the recall of the cross-attention network. A shortlist of the best 10\% of the ranking costs a tenth of docking the library, so that the shortlist saves a factor of ten in the docking of the library at the recall of Table~\ref{tab:known}. The docked training sample of 40{,}000 ligands per pocket is \trainShareLib{}\% of the library, so the first screen of the library, with the sample and the shortlist docked, costs \shareLabelShortlist{}\% of docking the library, a saving of a factor of \factorLabelShortlist{}. Every further library that the networks screen saves the factor of ten.

The networks serve the pockets of their training data. They were trained on a docked sample of 40{,}000 ligands per pocket and rank new ligands against these pockets, which is the setting of a screen with a docked sample for every pocket of interest. For a pocket of a protein family outside the training data the networks have no target-specific knowledge to apply: in a cross-validation over protein families the pocket state adds no recall for the held-out family (Supporting Information). Such a pocket needs docking scores of its own before this organization can screen it.
